% ============================================================
%  Modul 12 — Penjadwalan Proses Berbasis Prioritas
% ============================================================

\chapter{Penjadwalan Proses Berbasis Prioritas}
\label{chap:modul-12}

% --- 1. Tujuan dan Output Praktikum ---
\section{Tujuan dan Output Praktikum}

\subsection{Tujuan Praktikum}
Setelah menyelesaikan praktikum ini, mahasiswa diharapkan mampu:
\begin{enumerate}
	\item Memahami konsep \textit{Priority Scheduling} sebagai evolusi dari algoritma \textit{Round-Robin} yang telah dianalisis pada Modul~11.
	\item Memodifikasi fungsi \texttt{next\_thread\_to\_run()} di \texttt{src/thread.c} agar memilih \textit{thread} dengan nilai prioritas tertinggi dari \textit{ready queue}, menggantikan logika FIFO yang ada.
	\item Memverifikasi bahwa \textit{thread} berprioritas tinggi dieksekusi lebih dahulu melalui pengujian di QEMU dengan \textit{flag} \texttt{PRIORITY\_SCHEDULING}.
	\item Menganalisis fenomena \textit{starvation} yang muncul sebagai konsekuensi dari \textit{Priority Scheduling} murni dan memahami solusi konseptualnya.
\end{enumerate}

\subsection{Output Praktikum}
Pada akhir praktikum ini, mahasiswa diharapkan menghasilkan:
\begin{itemize}
	\item Implementasi fungsional \textit{Priority Scheduling} pada fungsi \texttt{next\_thread\_to\_run()} yang terbukti mampu memilih \textit{thread} berprioritas tertinggi dari \textit{ready queue}.
	\item Tangkapan layar (\textit{screenshot}) keluaran QEMU yang menunjukkan \textit{thread} \texttt{High} (prioritas 50) dieksekusi terlebih dahulu selama 5 tick, kemudian \texttt{Med} (prioritas 31) mendominasi CPU sepenuhnya, sementara \texttt{Low} (prioritas 10) mengalami \textit{starvation} dan tidak pernah mendapat giliran (selama \texttt{Med} masih ada di antrian).
	\item Laporan perbandingan yang menjabarkan perbedaan perilaku \textit{Round-Robin} (Modul~11) vs \textit{Priority Scheduling} (Modul~12) berdasarkan keluaran QEMU.
\end{itemize}

% --- 2. Dasar Teori ---
\newpage
\section{Dasar Teori}

\subsection{Priority Scheduling}
\textit{Priority Scheduling} adalah algoritma penjadwalan CPU yang menetapkan nilai prioritas numerik pada setiap \textit{thread} \cite{Silberschatz2018}. Ketika \textit{scheduler} harus memilih \textit{thread} berikutnya untuk dieksekusi, ia akan \textbf{memindai seluruh} \textit{ready queue} dan memilih \textit{thread} dengan nilai prioritas tertinggi, bukan yang paling depan di antrian.

Analoginya seperti ruang gawat darurat di rumah sakit. Meskipun ada antrean pasien yang datang lebih dulu, pasien dengan kondisi paling kritis (prioritas tertinggi) akan ditangani terlebih dahulu, tanpa memperhatikan urutan kedatangan.

Pada \textit{starter code} OS, prioritas didefinisikan melalui tiga konstanta di \texttt{src/thread.h}:

\begin{lstlisting}[language=C, caption={Konstanta prioritas pada thread.h}]
#define PRI_MIN     0    /* Prioritas terendah */
#define PRI_DEFAULT 31   /* Prioritas default */
#define PRI_MAX     63   /* Prioritas tertinggi */
\end{lstlisting}

Field \texttt{priority} pada \texttt{struct thread} sudah tersedia dan diinisialisasi sejak Modul~8, namun belum pernah dimanfaatkan oleh \textit{scheduler}. Pada modul ini, field tersebut akan menjadi kunci utama dalam menentukan giliran CPU.

\subsection{Perbedaan Mendasar dengan Round-Robin}
Perbedaan antara kedua algoritma terletak pada satu titik kritis: \textbf{bagaimana} \texttt{next\_thread\_to\_run()} memilih thread berikutnya dari antrian.

\begin{itemize}
	\item \textbf{Round-Robin}: Selalu mengambil thread dari posisi depan antrian (\texttt{ready\_head}). Tidak ada pemindaian, tidak ada perbandingan. Kompleksitas: $O(1)$.
	\item \textbf{Priority Scheduling}: Memindai \textit{seluruh} isi antrian untuk menemukan thread dengan \texttt{priority} tertinggi, mengeluarkannya dari posisinya (bukan dari depan), lalu menggeser elemen sisanya. Kompleksitas: $O(n)$, di mana $n$ adalah jumlah thread di antrian.
\end{itemize}

Berikut ilustrasi perbedaan keduanya pada antrian yang sama:

\begin{lstlisting}[caption={Ilustrasi perbedaan pemilihan thread}]
Ready Queue: [Low(10)] -> [Med(31)] -> [High(50)]

Round-Robin memilih : Low  (posisi depan, FIFO)
Priority memilih    : High (nilai priority tertinggi)
\end{lstlisting}

\begin{table}[h]
	\centering
	\caption{Perbandingan Round-Robin dan Priority Scheduling}
	\label{tab:m12-perbandingan}
	\begin{tabulary}{\textwidth}{|L|L|L|}
		\hline
		\textbf{Aspek} & \textbf{Round-Robin (Modul 11)} & \textbf{Priority Scheduling (Modul 12)} \\ \hline
		Prinsip pemilihan
		& FIFO: \textit{thread} yang lebih dulu masuk antrian, lebih dulu dieksekusi.
		& \textit{Thread} dengan nilai prioritas tertinggi selalu dipilih terlebih dahulu. \\ \hline
		Fairness
		& Tinggi: setiap \textit{thread} mendapat porsi CPU yang setara.
		& Rendah: \textit{thread} berprioritas rendah bisa tidak pernah mendapat CPU. \\ \hline
		Responsiveness
		& Rendah: \textit{thread} penting harus menunggu giliran yang sama.
		& Tinggi: \textit{thread} kritis langsung dieksekusi. \\ \hline
		Kompleksitas
		& $O(1)$: langsung ambil dari depan antrian.
		& $O(n)$: harus memindai seluruh antrian. \\ \hline
		Risiko utama
		& Tidak ada \textit{starvation}, tetapi tugas penting mungkin lambat ditangani.
		& \textit{Starvation} pada \textit{thread} berprioritas rendah. \\ \hline
	\end{tabulary}
\end{table}

\subsection{Starvation (Kelaparan)}
\textit{Starvation} adalah kondisi di mana sebuah \textit{thread} berprioritas rendah \textbf{tidak pernah} mendapat giliran CPU karena selalu ada \textit{thread} berprioritas lebih tinggi yang siap dieksekusi di \textit{ready queue}.

Pada skenario praktikum ini, \textit{starvation} total tidak terjadi karena \texttt{High} memiliki batas \texttt{max\_ticks = 5} dan akan \textit{exit} setelah 5 iterasi, sehingga akhirnya \texttt{Med} dan \texttt{Low} mendapat giliran. Namun, jika \texttt{High} berjalan tanpa batas, \texttt{Low} akan mengalami starvation.

Berikut skenario terjadinya \textit{starvation} langkah demi langkah:

\begin{table}[h]
	\centering
	\caption{Skenario starvation pada Priority Scheduling}
	\begin{tabulary}{\textwidth}{|c|l|L|}
		\hline
		\textbf{Langkah} & \textbf{Thread Aktif} & \textbf{Yang Terjadi} \\ \hline
		1 & High (pri=50) & Dipilih scheduler, berjalan, memanggil \texttt{yield()} \\ \hline
		2 & --- & Scheduler memindai antrian: Low(10), Med(31), High(50) \\ \hline
		3 & High (pri=50) & Dipilih lagi karena prioritas tertinggi \\ \hline
		4 & --- & High terus dipilih berulang kali \\ \hline
		$\infty$ & Low (pri=10) & \textbf{Tidak pernah mendapat giliran} = \textit{starvation} \\ \hline
	\end{tabulary}
\end{table}

\subsection{Solusi: Aging}
Teknik \textit{aging} adalah solusi klasik untuk mengatasi \textit{starvation}. Prinsipnya sederhana: setiap kali \textit{scheduler} dipanggil, prioritas semua thread yang \textbf{tidak} terpilih (masih menunggu di antrian) dinaikkan sebesar 1. Dengan demikian, thread berprioritas rendah yang sudah lama menunggu pada akhirnya akan memiliki prioritas cukup tinggi untuk mendapat giliran.

\begin{lstlisting}[caption={Pseudocode aging}]
function next_thread_to_run():
    chosen = thread dengan prioritas tertinggi

    for setiap thread LAIN di ready_queue:
        thread.priority += 1  /* Naikkan 1 setiap kali tidak terpilih */

    return chosen
\end{lstlisting}

Tabel \ref{tab:m12-dasar-teori} merangkum konsep-konsep penting pada modul ini.

\begin{table}[h]
	\caption{Tabel Ringkasan Dasar Teori Modul 12}
	\label{tab:m12-dasar-teori}
	\centering
	\begin{tabulary}{\textwidth}{|L|L|}
		\hline
		\textbf{Konsep} & \textbf{Penjelasan Singkat} \\ \hline
		Priority Scheduling
		& Algoritma penjadwalan yang memilih \textit{thread} berdasarkan nilai prioritas tertinggi di \textit{ready queue}. \\ \hline
		Starvation
		& Kondisi di mana \textit{thread} berprioritas rendah tidak pernah mendapat giliran CPU karena selalu dikalahkan oleh thread berprioritas lebih tinggi. \\ \hline
		Aging
		& Teknik pencegahan \textit{starvation} dengan menaikkan prioritas thread yang sudah lama menunggu. \\ \hline
		Pemindaian Antrian
		& Proses iterasi seluruh isi \textit{ready queue} untuk menemukan thread dengan \texttt{priority} tertinggi; kompleksitas $O(n)$. \\ \hline
	\end{tabulary}
\end{table}

% --- 3. Alat dan Bahan ---
\newpage
\section{Alat dan Bahan}
\label{sec:m12-alat-bahan}

\subsection{Perangkat Keras}
\begin{itemize}
	\item PC atau Laptop dengan arsitektur prosesor basis \textit{host} x86/x64 yang mendukung kelanjutan kompilasi emulator.
	\item Memori minimal 4GB RAM ruang eksekusi.
\end{itemize}

\subsection{Perangkat Lunak (Bahan Praktikum)}
\begin{itemize}
	\item OS \textit{Host} utama (Windows dengan WSL/GitBash, Linux, atau macOS).
	\item Text Editor (misalnya Visual Studio Code) untuk menyunting berkas kode bahasa C dan \textit{Assembly}.
	\item Instalasi Docker Engine (\texttt{docker.io}) untuk menyalakan lingkungan isolasi kompilasi.
	\item Emulator Hardware x86: \texttt{qemu-system-x86} (atau instalasi QEMU setara) yang akan mengeksekusi binari hasil OS.
	\item Direktori \textit{Starter Code} OS 32-bit dengan implementasi Modul~8 yang sudah berfungsi (berisi direktori \texttt{src} dengan kerangka lengkap \texttt{kernel.c}, \texttt{thread.h}, \texttt{thread.c}, \texttt{switch.asm}, dsb).
\end{itemize}

% --- 4. Langkah-Langkah Praktikum ---
\newpage
\section{Langkah-Langkah Praktikum}

\subsection{Persiapan Lingkungan}
\label{sec:m12-persiapan}

Pada praktikum ini, Anda akan memodifikasi fungsi \texttt{next\_thread\_to\_run()} di \texttt{src/thread.c} agar mengimplementasikan \textit{Priority Scheduling}. File yang dimodifikasi hanya \texttt{src/thread.c}. Pastikan Anda menggunakan direktori kerja yang sama dengan Modul~11.

\textit{Catatan: Pastikan Anda masih berada di dalam lingkungan Docker yang digunakan pada Modul 11, atau jalankan kembali \texttt{./run-docker.sh} jika sudah keluar.}

\begin{enumerate}
	\item Buka terminal, pastikan Anda berada di dalam direktori \textit{source code} OS (direktori kerja yang sama dengan Modul~11).
	\item Bersihkan hasil \textit{build} sebelumnya agar kompilasi berikutnya menggunakan kode terbaru yang akan Anda tulis.
	      \begin{lstlisting}[language=bash, caption={Bersihkan build sebelumnya}]
make clean
    \end{lstlisting}
\end{enumerate}

% -----------------------------------------------------------------
\subsection{Sesi 1: Analisis Kode Scheduler Round-Robin yang Akan Dirombak}
\label{sec:m12-sesi-1}

Sebelum menulis kode baru, pahami secara menyeluruh kode \textit{Round-Robin} yang saat ini ada. Buka file \texttt{src/thread.c} menggunakan teks editor Anda dan cari fungsi \texttt{next\_thread\_to\_run()}.

\begin{enumerate}
	\item Perhatikan implementasi \textit{Round-Robin} (FIFO) yang Anda tulis pada Modul~8:
	      \begin{lstlisting}[language=C, caption={Implementasi Round-Robin saat ini di next\_thread\_to\_run()}]
struct thread *next_thread_to_run(void) {
    if (ready_count == 0)
        return initial_thread;
    struct thread *next = ready_queue[ready_head];
    ready_head = (ready_head + 1) % (MAX_THREADS * 2);
    ready_count--;
    return next;
}
    \end{lstlisting}

	\item Identifikasi kelemahan kode di atas: baris \texttt{ready\_queue[ready\_head]} selalu mengambil \textit{thread} dari posisi depan antrian tanpa memeriksa field \texttt{priority}. Akibatnya, \textit{thread} berprioritas tinggi diperlakukan sama dengan \textit{thread} berprioritas rendah, sebagaimana telah dibuktikan pada keluaran QEMU di Modul~11.

	\item Perhatikan juga variabel-variabel antrian yang akan Anda gunakan kembali:
	      \begin{lstlisting}[language=C, caption={Variabel antrian thread di thread.c (sudah disediakan)}]
static struct thread *ready_queue[MAX_THREADS * 2];
static uint32_t ready_head  = 0;
static uint32_t ready_tail  = 0;
static uint32_t ready_count = 0;
    \end{lstlisting}
	      Antrian ini berupa \textit{circular array} dengan kapasitas \texttt{MAX\_THREADS * 2} (16 slot). Variabel \texttt{ready\_head} menunjuk ke elemen terdepan, \texttt{ready\_tail} menunjuk ke posisi kosong berikutnya, dan \texttt{ready\_count} menghitung jumlah thread yang saat ini ada di antrian.
\end{enumerate}

% -----------------------------------------------------------------
\subsection{Sesi 2: Implementasi Priority Scheduling}
\label{sec:m12-sesi-2}

Modifikasi fungsi \texttt{next\_thread\_to\_run()} agar memilih \textit{thread} dengan prioritas tertinggi dari \textit{ready queue}, menggantikan logika FIFO. Masih di dalam file \texttt{src/thread.c}:

\begin{enumerate}
	\item \textbf{Implementasi \texttt{next\_thread\_to\_run()} (Priority Scheduling)} \\
	      Ganti seluruh isi fungsi \texttt{next\_thread\_to\_run()} dengan kode berikut. Logikanya terdiri dari tiga tahap yang ditandai dengan komentar bernomor:

	      \begin{lstlisting}[language=C, caption={Kode implementasi Priority Scheduling pada next\_thread\_to\_run()}]
struct thread *next_thread_to_run(void) {
    if (ready_count == 0)
        return initial_thread;

    /* 1. Cari indeks thread dengan prioritas tertinggi */
    uint32_t best_idx  = ready_head;
    int      best_prio = ready_queue[ready_head]->priority;

    for (uint32_t i = 1; i < ready_count; i++) {
        uint32_t idx = (ready_head + i) % (MAX_THREADS * 2);
        if (ready_queue[idx]->priority > best_prio) {
            best_prio = ready_queue[idx]->priority;
            best_idx  = idx;
        }
    }

    /* 2. Ambil thread terpilih dari antrian */
    struct thread *next = ready_queue[best_idx];

    /* 3. Geser elemen setelah posisi terpilih untuk menutup celah */
    uint32_t i = best_idx;
    while (1) {
        uint32_t next_i = (i + 1) % (MAX_THREADS * 2);
        if (next_i == ready_tail) break;
        ready_queue[i] = ready_queue[next_i];
        i = next_i;
    }
    ready_tail  = (ready_tail == 0)
                      ? (MAX_THREADS * 2 - 1)
                      : ready_tail - 1;
    ready_count--;

    return next;
}
    \end{lstlisting}

	      Penjelasan setiap tahap:
	      \begin{itemize}
		      \item \textbf{Tahap 1 (Pemindaian):} Loop menelusuri seluruh isi antrian mulai dari \texttt{ready\_head}. Untuk setiap elemen, bandingkan \texttt{priority}-nya dengan \texttt{best\_prio}. Jika lebih tinggi, catat indeks aktualnya di \texttt{best\_idx} dan perbarui \texttt{best\_prio}. Setelah loop selesai, \texttt{best\_idx} menunjuk ke posisi thread dengan prioritas tertinggi.

		      \item \textbf{Tahap 2 (Pengambilan):} Simpan pointer thread terpilih dari \texttt{ready\_queue[best\_idx]} ke variabel \texttt{next}.

		      \item \textbf{Tahap 3 (Penggeseran):} Karena thread yang diambil mungkin bukan dari posisi depan antrian (berbeda dengan FIFO), elemen-elemen setelahnya harus digeser maju satu posisi untuk menutup celah. Ini dilakukan oleh loop \texttt{while} yang bergerak dari \texttt{best\_idx} hingga \texttt{ready\_tail}. Terakhir, mundurkan \texttt{ready\_tail} dan kurangi \texttt{ready\_count}.
	      \end{itemize}

	      \textit{Catatan penting: Jika muncul error \texttt{implicit declaration of function} saat compile, pastikan Anda tidak menghapus baris \texttt{\#include "thread.h"} dan \texttt{\#include "types.h"} di awal file.}

	      \textit{Jika muncul warning \texttt{comparison of unsigned expression < 0 is always true}, pastikan tipe data \texttt{ready\_count} adalah \texttt{uint32\_t} (bukan \texttt{int}). Cek deklarasi variabel di bagian atas file \texttt{thread.c}.}

	\item \textbf{Simpan} file \texttt{thread.c}. Saat ini \textit{scheduler} Anda sudah berubah dari \textit{Round-Robin} menjadi \textit{Priority Scheduling}.
\end{enumerate}

% -----------------------------------------------------------------
\subsection{Sesi 3: Kompilasi dan Verifikasi Sistem OS}
\label{sec:m12-sesi-3}

\begin{enumerate}
	\item Lakukan kompilasi sistem operasi Anda dengan \textit{flag} \texttt{PRIORITY\_SCHEDULING}.
	      \begin{lstlisting}[language=bash, caption={Compile dan jalankan dengan Priority Scheduling}]
# Pastikan telah menjalankan make clean sebelumnya
make DEMO=-DPRIORITY_SCHEDULING
make run
    \end{lstlisting}

	      \textit{Jika muncul error \texttt{linker error: undefined reference to `next\_thread\_to\_run'}, berarti ada kesalahan sintaks pada fungsi Anda (misalnya kurung kurawal tidak tertutup). Periksa kembali kode pada Sesi 2.}

	      \textit{Jika compile berhasil tetapi QEMU langsung menampilkan \texttt{KERNEL PANIC}, kemungkinan \texttt{next\_thread\_to\_run()} mengembalikan \texttt{NULL}. Pastikan blok \texttt{if (ready\_count == 0) return initial\_thread;} masih ada di awal fungsi.}

	\item \textbf{Verifikasi Hasil:} Amati keluaran (\textit{output}) di layar QEMU. Bandingkan dengan perilaku sebelum modifikasi.
	      \begin{itemize}
		      \item \textbf{Sebelumnya (Modul 11):} Ketiga \textit{thread} bergantian secara merata tanpa memperhatikan prioritas (\textit{Round-Robin}).
		      \item \textbf{Sekarang (Modul 12):} Thread \texttt{High} (prioritas 50) mendominasi CPU selama 5 tick penuh. Setelah \texttt{High} \textit{exit}, \texttt{Med} (prioritas 31) mendominasi CPU sepenuhnya, sementara \texttt{Low} (prioritas 10) \textbf{tidak pernah mendapat giliran} (\textit{starvation}).
	      \end{itemize}

	      \begin{lstlisting}[caption={Ekspektasi keluaran QEMU setelah Priority Scheduling}]
===================================
  IF-OS  |  Praktikum Sistem Operasi
===================================

Inisialisasi sistem thread...
Demo Minggu 6: Priority Scheduling
(Jika masih Round-Robin, next_thread_to_run() belum dimodifikasi)
(High=50 jalan 5 tick lalu exit, Med=31 mengambil alih, Low=10 terakhir)

[High] tick 0
[High] tick 1
[High] tick 2
[High] tick 3
[High] tick 4
[High] selesai (5 tick)
[Med] tick 0
[Med] tick 1
[Med] tick 2
[Med] tick 3
...            <- Med berjalan terus, Low tidak pernah muncul
    \end{lstlisting}

	      Meskipun \textit{thread} \texttt{Low} memanggil \texttt{yield()} dan masuk kembali ke antrian, \textit{scheduler} selalu memilih \textit{thread} berprioritas tertinggi. Karena \texttt{Med} (31) selalu memiliki prioritas lebih tinggi daripada \texttt{Low} (10), \texttt{Med} akan terus dipilih dan \texttt{Low} \textbf{tidak pernah mendapat giliran sama sekali}.

	      Fenomena ini adalah contoh nyata dari \textbf{\textit{starvation}} yang telah dibahas pada Dasar Teori. Pada skenario ini, \texttt{High} tidak menyebabkan \textit{starvation} permanen karena memiliki batas \texttt{max\_ticks = 5}, tetapi \texttt{Med} (yang berjalan tanpa batas) secara efektif membuat \texttt{Low} kelaparan selamanya. Untuk mengatasi ini, diperlukan mekanisme \textit{aging} sebagaimana dijelaskan pada Dasar Teori, namun implementasinya berada di luar cakupan modul ini.




	\item Ambil tangkapan layar (\textit{screenshot}) hasil eksekusi untuk membuktikan bahwa implementasi \textit{Priority Scheduling} telah berjalan sesuai ekspektasi.
\end{enumerate}

\begin{figure}[h]
	\centering
	% TODO: ganti dengan screenshot hasil QEMU Priority Scheduling
	\includegraphics[width=0.85\textwidth]{Figure/12/1.png}
	\caption{Keluaran (\textit{output}) QEMU setelah implementasi Priority Scheduling. Thread High (prioritas 50) dieksekusi pertama kali selama 5 tick penuh, kemudian Med (prioritas 31) mendominasi CPU sepenuhnya sementara Low (prioritas 10) mengalami \textit{starvation}.}
	\label{fig:m12-sesi3}
\end{figure}
